\subsection{Absolute convergence at the boundary}
\label{audit:poly:convergence}

\paragraph{Source.}
\texttt{multiple-polylogarithms-introduction.tex}, the theorem titled
``Domain of absolute convergence'' (source lines 231--235).
The condition printed there, $\Re s_j\geq1$, $|z_j|\leq1$, and
$(s_1,z_1)\neq(1,1)$, does not imply absolute convergence.
For example, $\Li_1(i)$ satisfies every stated hypothesis, while the
absolute-value series is $\sum_{n\geq1}1/n$.

\paragraph{Replacement.}
For $\Re s_j\geq1$ and $|z_j|\leq1$, a sufficient condition for
absolute convergence is
\[
 |z_1|<1\qquad\text{or}\qquad\Re s_1>1.
\]
Indeed, the absolute value of the inner nested sum is at most
$H_{n_1-1}^{k-1}/(k-1)!$, so the full series is dominated by
\[
 \frac1{(k-1)!}\sum_{n\geq1}
    \frac{|z_1|^n(1+\log n)^{k-1}}{n^{\Re s_1}}.
\]
Conditional convergence at boundary roots of unity is a separate issue.
The usual positive-integer MZV admissibility condition $s_1\geq2$ remains
correct. Also, the definition of $S_p$ in
\texttt{gaussian-multiple-polylog-depth.tex} should start at $p\geq1$:
its printed allowance $p=0$ gives terms $(-1)^nH_n$ that do not tend to
zero. An Abel-regularized value would need its own explicit definition.

\subsection{The order of the iterated-integral letters}
\label{audit:poly:integral-order}

\paragraph{Source.}
\texttt{gaussian-eisenstein-double-polylogs.tex},
\texttt{eq:G-def} and \texttt{eq:goncharov-dbl}; the same mismatch appears
in \texttt{multiple-polylogarithms-introduction.tex}, lines 276--282.
The displayed definition attaches the letter $a_i$ to $t_i$ on the simplex
$0<t_1<\cdots<t_n<1$, but the subsequent polylogarithm dictionary uses the
opposite convention. In particular, the displayed integral for
$G(0,z^{-1};1)$ contains $\int_0^{t_2}dt_1/t_1$, which diverges, although
$-\Li_2(z)$ is finite, for example at $z=1/2$.

\paragraph{Replacement.}
Keep the stated dictionary and define the leftmost letter to be outermost:
\begin{align*}
 G(a_1,\ldots,a_n;z)
   &=\int_0^z\frac{G(a_2,\ldots,a_n;t)}{t-a_1}\,dt,\\
 G(a_1,\ldots,a_n;1)
   &=\int_{0<t_n<\cdots<t_1<1}
            \prod_{j=1}^n\frac{dt_j}{t_j-a_j}.
\end{align*}
For convergent words this gives
\[
 \Li_n(z)=-G(0^{n-1},z^{-1};1),\qquad
 \Li_{a,b}(u,v)=G(0^{a-1},u^{-1},0^{b-1},(uv)^{-1};1).
\]
Alternatively, one may retain the increasing-index simplex and reverse
every letter word in the dictionary. The two choices must not be mixed.
Endpoint regularization, when needed, should be stated separately.

\subsection{An extra term in the modified polylogarithm}
\label{audit:poly:modified}

\paragraph{Source.}
\path{ladders-as-bloch-elements__ce20ea60a460.tex},
\texttt{eq:Pm}, and the reported complex-embedding value of $P_4$.
The upper endpoint of the sum is printed as $m$. For $m\geq2$ the correct
modified polylogarithm is
\begin{equation}\label{audit:poly:correct-P}
 P_m(z)=\operatorname{Re}_m\!\left(
       \sum_{j=0}^{m-1}\frac{2^jB_j}{j!}
          \log^j|z|\,\Li_{m-j}(z)\right),
 \quad
 \operatorname{Re}_m=
 \begin{cases}\Re,&m\ \text{odd},\\\Im,&m\ \text{even}.
 \end{cases}
\end{equation}
Here $B_1=-1/2$. This agrees with the standard $P_m$ definition in the
PARI/GP documentation, \texttt{polylog} with flag $3$
\cite{audit-pari-polylog}; its additional real multiple of
$B_m\log^m|z|$ vanishes after $\operatorname{Re}_m$ for $m\geq2$.

The printed expression adds, at even $m$, the generally nonzero term
\begin{equation}\label{audit:poly:P-error}
 \frac{2^mB_m}{m!}\log^m|z|\,
             \Im\!\left(\frac{z}{1-z}\right).
\end{equation}
For example, at $m=2$, $z=i/2$, the excess is exactly
$2\log^2(2)/15$. Verification only at odd weights or at $z=i$ cannot
detect this error, because respectively $B_m=0$ or $\log|z|=0$.

The error also affects a numerical conclusion. For the lower nonreal
root of $z^3+z-1=0$,
\[
z_-=-0.34116390191400966\ldots-1.16154139999725194\ldots i,
\]
recomputation gives
\begin{align*}
 P_4(z_-)&=-0.916010467826680838722383200166\ldots,\\
 P_4^{\mathrm{printed}}(z_-)&=-0.915999527027516872882916723105\ldots.
\end{align*}
The latter agrees with the report's $-0.91599\ldots$.
Thus the definition and the complex-embedding computation both require
correction; the odd-weight formulas are not affected by this particular
endpoint error.

\subsection{The asserted binomial dimensions are not dimensions of depth}
\label{audit:poly:depth-dimensions}

\paragraph{Source.}
\texttt{gaussian-eisenstein-double-polylogs.tex},
\texttt{thm:deligne}, \texttt{eq:binomial-dim}, and
\texttt{cor:depth3}. The formula
$\dim\operatorname{gr}^D_p(\text{weight }n)=\binom np$
is false for the ordinary cyclotomic depth filtration used in the cited
literature.

A direct contradiction already occurs at $n=6$, $p=1$. All level-four
weight-six single polylogarithms lie in the span of $\zeta(6)$ and
$i\beta(6)$, since
\[
 \Li_6(-1)=-\frac{31}{32}\zeta(6),\qquad
 \Li_6(\pm i)=-\frac{31}{2048}\zeta(6)\pm i\beta(6).
\]
The depth-one space therefore has dimension at most $2$, whereas the
printed formula gives $6$. Likewise, the depth-zero component of positive
weight is zero, not $\binom n0=1$.
These contradictions use explicit identities and do not depend on
unproved numerical-period independence.

Glanois defines depth by the number of nonzero iterated-integral letters
and explicitly warns that it is a filtration rather than a grading on
$\mathcal H$ \cite[\S3.1]{audit-glanois-2016}.
A Hilbert series in weight alone does not determine these filtered
pieces. Moreover, even a correct positive dimension for an ambient
motivic depth-three space would not imply that three particular
$\Li_{a,b,c}(i,1,1)$ have nonzero independent images in it.
An explicit computation of their classes is required.
The period conjecture does not repair either of these logical gaps.

\paragraph{Replacement.}
Withdraw \texttt{eq:binomial-dim} and the stated proof of
\texttt{cor:depth3}. Retain the verified shuffle relations as reductions
modulo their displayed product terms. The minimal depth of the three
specified numerical survivors remains undetermined by the arguments in
that draft. A motivic claim should name the precise filtration and supply
the relevant coaction or basis computation before making a corresponding
conditional numerical claim.

\subsection{Rank seven does not mean seven individual triples collapse}
\label{audit:poly:shuffle-rank}

\paragraph{Source.}
\texttt{gaussian-eisenstein-double-polylogs.tex},
\texttt{prop:7of10} and its preceding rank calculation.
There are ten compositions $(a,b,c)$ of weight $6$ with positive entries.
The draft obtains rank $7$ from shuffle products and concludes that seven
of the ten triples are themselves products. The correct conclusion is
that seven independent \emph{linear combinations} are products.

We recomputed the $13$ formal shuffle-product rows exactly: the ten
products of a single word and a double word, together with the three
unordered triple products of single words. Their rational rank is indeed
$7$. However, none of the ten individual coordinate vectors belongs to
the row space of these products.
This can be checked using the following three integer linear functionals,
in lexicographic order of the compositions
$(1,1,4),(1,2,3),\ldots,(4,1,1)$:
\[
 \begin{pmatrix}
 -1&2&-2&1&0&0&-1&1&0&0\\
 -3&6&-6&3&-1&4&-4&0&1&0\\
 -6&12&-12&6&-3&12&-9&0&0&1
 \end{pmatrix}.
\]
Each row annihilates every shuffle-product row, the three rows are
independent, and each of the ten columns is nonzero.
Thus every individual word is detected by at least one annihilating
functional. The full matrix and exact checks are generated by
\texttt{code/verify\_shuffle\_audit.py} and recorded in
\texttt{data/gaussian\_shuffle\_audit.json}.

\paragraph{Replacement.}
The formal weight-six shuffle quotient modulo products is
three-dimensional, and seven of its ten coordinate variables can be
eliminated \emph{in terms of three survivors and products}.
This formal calculation neither proves nor disproves additional
relations after evaluation at $i$.

The same distinction applies to \texttt{thm:gauss-5} and
\texttt{thm:eis-1}: a corank computed from a specified relation matrix is
an upper bound on the dimension of the corresponding numerical span,
unless completeness of those relations is independently established.

\subsection{Two further logical boundary corrections}
\label{audit:poly:other}

In \texttt{gaussian-multiple-polylog-depth.tex}, the discussion of the
weight-four symmetric triple sum compares it to
$\zeta(2,1,1)+\zeta(1,2,1)+\zeta(1,1,2)$ as though this were a convergent
uncolored MZV combination. The last two defining sums diverge. In each,
the inner sum is bounded below by a positive constant for all sufficiently
large outer indices, and the outer exponent is $1$; comparison with the
harmonic series proves divergence. The comparison should be removed,
or replaced by an explicitly specified regularized identity.

The same draft infers at least $50$ independent new directions from
$50$ individual failures of membership in a fixed span. Even exact
non-membership of each value would not justify that count: all $50$
values could represent one common nonzero class in the quotient.
Numerical failures of membership provide still less. The appropriate
replacement is a list of tested candidate values with no claimed lower
bound on the dimension of their joint span.
